% =====================================================================================
% ANEXO A - Evidencias de existencia de las fuentes referidas
% Documento independiente del reporte de estadía (la plantilla oficial no se modifica).
% Compilar con: pdflatex anexo_evidencias.tex
% =====================================================================================
\documentclass[12pt,letterpaper]{article}

\usepackage[left=2.5cm,top=2.0cm,right=2.5cm,bottom=3.0cm]{geometry}
\usepackage[utf8]{inputenc}
\usepackage[T1]{fontenc}
\usepackage[spanish,mexico]{babel}
\usepackage{tabularx}
\usepackage{booktabs}
\usepackage{float}
\usepackage{enumitem}
\usepackage{fancyhdr}
\usepackage[hidelinks]{hyperref}
\usepackage{titlesec}
\usepackage{setspace}

\setlength{\parskip}{0.6em}
\setlength{\parindent}{0pt}
\onehalfspacing

\pagestyle{fancy}
\fancyhf{}
\fancyhead[L]{\footnotesize Anexo A. Evidencias de existencia de las fuentes referidas}
\fancyhead[R]{\footnotesize\thepage}
\renewcommand{\headrulewidth}{0.4pt}

\begin{document}

\begin{center}
  {\large\bfseries Universidad Politécnica de Victoria}\\[2pt]
  {\normalsize Ingeniería en Tecnologías de la Información e Innovación Digital}\\[10pt]
  {\Large\bfseries Anexo A}\\[3pt]
  {\large\bfseries Evidencias de existencia de las fuentes referidas}\\[10pt]
  {\normalsize Joshua André Alvarado Tovar \quad --- \quad Matrícula 2330242}\\[2pt]
  {\footnotesize Estadía en el Instituto Registral y Catastral del Estado de Tamaulipas}\\[2pt]
  {\footnotesize Cuatrimestre agosto--diciembre de 2026}
\end{center}

\vspace{6pt}
\hrule
\vspace{10pt}

\section*{Propósito y método de verificación}

Este anexo agrupa la evidencia de existencia de las veintitrés fuentes citadas en el reporte de estadía, en atención
al requisito de acreditar que cada referencia corresponde a un documento real y localizable. Se presenta como archivo
independiente, de modo que la estructura del reporte permanezca conforme a la plantilla institucional.

Cada fuente se consigna en una ficha que contiene su referencia completa, el tipo de documento de que se trata, el
identificador que la hace irrepetible ---norma, ISBN, fecha de publicación oficial o dominio institucional---, la
localización en que puede recuperarse y la constancia de la verificación practicada.

El criterio seguido para la verificación fue el siguiente. Tratándose de normas internacionales, se confirmó el
número de registro en el catálogo del organismo emisor y, cuando existe, en los catálogos nacionales que adoptan la
norma. Tratándose de libros, se confirmó el ISBN de la edición citada y su presencia en el catálogo del editor o en
catálogos bibliotecarios. Tratándose de legislación, se confirmó la fecha de publicación en el órgano de difusión
oficial correspondiente y la disponibilidad del texto vigente en el repositorio del poder legislativo que lo expidió.
Tratándose de documentación técnica y de sitios de producto, se confirmó que el recurso se publica en el dominio
oficial del organismo o del proyecto que lo sostiene.

La verificación se practicó el 29 de octubre de 2026. Las fichas se numeran de A-01 a A-23 y se ordenan conforme a los
bloques temáticos del archivo de bibliografía del reporte; la clave de cita que encabeza cada ficha es la misma que se
emplea en el documento, de manera que toda cita del texto puede rastrearse hasta su evidencia.

\vspace{4pt}


\subsection*{Marco jurídico y organismo receptor}
\addcontentsline{toc}{subsection}{Marco jurídico y organismo receptor}

\begin{table}[H]
  \small
  \renewcommand{\arraystretch}{1.25}
  \begin{tabularx}{\textwidth}{@{}>{\raggedright\arraybackslash}p{3.0cm}>{\raggedright\arraybackslash}X@{}}
    \toprule
    \multicolumn{2}{@{}l}{\textbf{Ficha A-01} \quad\textnormal{\footnotesize clave de cita: \texttt{poe2022irct}}} \\
    \midrule
    \textbf{Referencia} & Gobierno del Estado de Tamaulipas (2022). \emph{Periódico Oficial del Estado de Tamaulipas}, Tomo CXLVII, número 116, 28 de septiembre de 2022. \\
    \textbf{Tipo de fuente} & Publicación oficial del Estado \\
    \textbf{Identificador} & Tomo CXLVII, núm.\ 116 (28-09-2022) \\
    \textbf{Localización} & \texttt{\footnotesize po.tamaulipas.gob.mx/wp-content/uploads/2022/09/\allowbreak cxlvii-116-280922F.pdf} \\
    \textbf{Constancia} & El ejemplar se localiza en el repositorio del Periódico Oficial del Estado y está indizado de manera independiente en la base legislativa vLex con el registro \emph{Periódico Oficial Legislativo --- Número 116 --- Tomo CXLVII del 28-09-2022}. \\
    \bottomrule
  \end{tabularx}
\end{table}

\begin{table}[H]
  \small
  \renewcommand{\arraystretch}{1.25}
  \begin{tabularx}{\textwidth}{@{}>{\raggedright\arraybackslash}p{3.0cm}>{\raggedright\arraybackslash}X@{}}
    \toprule
    \multicolumn{2}{@{}l}{\textbf{Ficha A-02} \quad\textnormal{\footnotesize clave de cita: \texttt{leyrpp2008}}} \\
    \midrule
    \textbf{Referencia} & H. Congreso del Estado de Tamaulipas (2008). \emph{Ley del Registro Público de la Propiedad Inmueble y del Comercio}. Última reforma del 6 de junio de 2023. \\
    \textbf{Tipo de fuente} & Legislación estatal vigente \\
    \textbf{Identificador} & P.O.\ anexo al núm.\ 27 (28-02-2008) \\
    \textbf{Localización} & \texttt{\footnotesize congresotamaulipas.gob.mx/Parlamentario/Archivos/Leyes/} \\
    \textbf{Constancia} & El texto vigente se publica en el apartado de leyes del H. Congreso del Estado de Tamaulipas y se reproduce además en el repositorio del Periódico Oficial del Estado. \\
    \bottomrule
  \end{tabularx}
\end{table}

\begin{table}[H]
  \small
  \renewcommand{\arraystretch}{1.25}
  \begin{tabularx}{\textwidth}{@{}>{\raggedright\arraybackslash}p{3.0cm}>{\raggedright\arraybackslash}X@{}}
    \toprule
    \multicolumn{2}{@{}l}{\textbf{Ficha A-03} \quad\textnormal{\footnotesize clave de cita: \texttt{leycatastro2001}}} \\
    \midrule
    \textbf{Referencia} & H. Congreso del Estado de Tamaulipas (2001). \emph{Ley de Catastro para el Estado de Tamaulipas}. Última reforma del 21 de diciembre de 2021. \\
    \textbf{Tipo de fuente} & Legislación estatal vigente \\
    \textbf{Identificador} & P.O.\ núm.\ 130 (30-10-2001) \\
    \textbf{Localización} & \texttt{\footnotesize congresotamaulipas.gob.mx/Parlamentario/Archivos/Leyes/} \\
    \textbf{Constancia} & El texto vigente se publica en el apartado de leyes del H. Congreso del Estado de Tamaulipas, que conserva asimismo el archivo de la reforma de diciembre de 2021. \\
    \bottomrule
  \end{tabularx}
\end{table}

\begin{table}[H]
  \small
  \renewcommand{\arraystretch}{1.25}
  \begin{tabularx}{\textwidth}{@{}>{\raggedright\arraybackslash}p{3.0cm}>{\raggedright\arraybackslash}X@{}}
    \toprule
    \multicolumn{2}{@{}l}{\textbf{Ficha A-04} \quad\textnormal{\footnotesize clave de cita: \texttt{lga2018}}} \\
    \midrule
    \textbf{Referencia} & Congreso General de los Estados Unidos Mexicanos (2018). \emph{Ley General de Archivos}. Diario Oficial de la Federación, 15 de junio de 2018. \\
    \textbf{Tipo de fuente} & Legislación federal vigente \\
    \textbf{Identificador} & D.O.F.\ 15-06-2018 \\
    \textbf{Localización} & \texttt{\footnotesize diputados.gob.mx/LeyesBiblio/pdf/LGA.pdf} \\
    \textbf{Constancia} & La Cámara de Diputados conserva el texto vigente en LeyesBiblio y el texto original de publicación en \texttt{LeyesBiblio/ref/lga/LGA\_orig\_15jun18.pdf}, junto con la relación histórica de reformas. \\
    \bottomrule
  \end{tabularx}
\end{table}

\begin{table}[H]
  \small
  \renewcommand{\arraystretch}{1.25}
  \begin{tabularx}{\textwidth}{@{}>{\raggedright\arraybackslash}p{3.0cm}>{\raggedright\arraybackslash}X@{}}
    \toprule
    \multicolumn{2}{@{}l}{\textbf{Ficha A-05} \quad\textnormal{\footnotesize clave de cita: \texttt{lgpdppso2017}}} \\
    \midrule
    \textbf{Referencia} & Congreso General de los Estados Unidos Mexicanos (2017). \emph{Ley General de Protección de Datos Personales en Posesión de Sujetos Obligados}. Diario Oficial de la Federación, 26 de enero de 2017. \\
    \textbf{Tipo de fuente} & Legislación federal vigente \\
    \textbf{Identificador} & D.O.F.\ 26-01-2017 \\
    \textbf{Localización} & \texttt{\footnotesize diputados.gob.mx/LeyesBiblio/pdf/LGPDPPSO.pdf} \\
    \textbf{Constancia} & La Cámara de Diputados publica el texto vigente en LeyesBiblio y mantiene la ficha histórica de la ley en \texttt{LeyesBiblio/abro/lgpdppso\_2017.htm}. \\
    \bottomrule
  \end{tabularx}
\end{table}

\begin{table}[H]
  \small
  \renewcommand{\arraystretch}{1.25}
  \begin{tabularx}{\textwidth}{@{}>{\raggedright\arraybackslash}p{3.0cm}>{\raggedright\arraybackslash}X@{}}
    \toprule
    \multicolumn{2}{@{}l}{\textbf{Ficha A-06} \quad\textnormal{\footnotesize clave de cita: \texttt{sgg2023irct}}} \\
    \midrule
    \textbf{Referencia} & Secretaría General de Gobierno del Estado de Tamaulipas (2023). \emph{Instituto Registral y Catastral}. \\
    \textbf{Tipo de fuente} & Sitio institucional oficial \\
    \textbf{Identificador} & Dominio de gobierno \texttt{tamaulipas.gob.mx} \\
    \textbf{Localización} & \texttt{\footnotesize tamaulipas.gob.mx/secretariageneral/\allowbreak instituto-registral-y-catastral/} \\
    \textbf{Constancia} & Página institucional del organismo dentro del portal de la Secretaría General de Gobierno. La existencia jurídica del Instituto queda acreditada de manera independiente por el decreto de creación referido en la ficha A-01. \\
    \bottomrule
  \end{tabularx}
\end{table}

\begin{table}[H]
  \small
  \renewcommand{\arraystretch}{1.25}
  \begin{tabularx}{\textwidth}{@{}>{\raggedright\arraybackslash}p{3.0cm}>{\raggedright\arraybackslash}X@{}}
    \toprule
    \multicolumn{2}{@{}l}{\textbf{Ficha A-07} \quad\textnormal{\footnotesize clave de cita: \texttt{transparenciairct}}} \\
    \midrule
    \textbf{Referencia} & Gobierno del Estado de Tamaulipas (2023). \emph{Instituto Registral y Catastral de Tamaulipas. Información pública, entidades}. \\
    \textbf{Tipo de fuente} & Portal oficial de transparencia \\
    \textbf{Identificador} & Dominio \texttt{transparencia.tamaulipas.gob.mx} \\
    \textbf{Localización} & \texttt{\footnotesize transparencia.tamaulipas.gob.mx/informacion-publica/\allowbreak entidades/instituto-registral-y-catastral-de-tamaulipas/} \\
    \textbf{Constancia} & El organismo figura como entidad sujeta a obligaciones de transparencia; el mismo portal publica su cuenta pública anual, lo que constituye evidencia documental adicional de su existencia y operación. \\
    \bottomrule
  \end{tabularx}
\end{table}

\begin{table}[H]
  \small
  \renewcommand{\arraystretch}{1.25}
  \begin{tabularx}{\textwidth}{@{}>{\raggedright\arraybackslash}p{3.0cm}>{\raggedright\arraybackslash}X@{}}
    \toprule
    \multicolumn{2}{@{}l}{\textbf{Ficha A-08} \quad\textnormal{\footnotesize clave de cita: \texttt{tamaulipas2025mod}}} \\
    \midrule
    \textbf{Referencia} & Gobierno del Estado de Tamaulipas (2025). \emph{Instituto Registral y Catastral de Tamaulipas fortalece la Contraloría Social y presenta avances del Proyecto de Modernización Integral 2025}. \\
    \textbf{Tipo de fuente} & Comunicado oficial de gobierno \\
    \textbf{Identificador} & Dominio de gobierno \texttt{tamaulipas.gob.mx} \\
    \textbf{Localización} & \texttt{\footnotesize tamaulipas.gob.mx/2025/11/instituto-registral-y-catastral-de-\allowbreak tamaulipas-fortalece-la-contraloria-social/} \\
    \textbf{Constancia} & Comunicado publicado en la sala de prensa del portal oficial del Gobierno del Estado, con fecha de publicación verificable en la propia dirección del recurso. \\
    \bottomrule
  \end{tabularx}
\end{table}

\begin{table}[H]
  \small
  \renewcommand{\arraystretch}{1.25}
  \begin{tabularx}{\textwidth}{@{}>{\raggedright\arraybackslash}p{3.0cm}>{\raggedright\arraybackslash}X@{}}
    \toprule
    \multicolumn{2}{@{}l}{\textbf{Ficha A-09} \quad\textnormal{\footnotesize clave de cita: \texttt{futuver2024}}} \\
    \midrule
    \textbf{Referencia} & Futuver Grupo (2024). \emph{Modernización Registral Electrónica: Registro Público, Estado de Tamaulipas, México}. \\
    \textbf{Tipo de fuente} & Documento corporativo del proveedor \\
    \textbf{Identificador} & Dominio corporativo \texttt{futuver.com} \\
    \textbf{Localización} & \texttt{\footnotesize futuver.com/modernizacion-registral-rpp-tamaulipas/} \\
    \textbf{Constancia} & Ficha de proyecto publicada por la empresa responsable de la implementación del sistema registral electrónico del Estado, en su sitio corporativo. \\
    \bottomrule
  \end{tabularx}
\end{table}

\subsection*{Normas y estándares}
\addcontentsline{toc}{subsection}{Normas y estándares}

\begin{table}[H]
  \small
  \renewcommand{\arraystretch}{1.25}
  \begin{tabularx}{\textwidth}{@{}>{\raggedright\arraybackslash}p{3.0cm}>{\raggedright\arraybackslash}X@{}}
    \toprule
    \multicolumn{2}{@{}l}{\textbf{Ficha A-10} \quad\textnormal{\footnotesize clave de cita: \texttt{iso15489}}} \\
    \midrule
    \textbf{Referencia} & International Organization for Standardization (2016). \emph{ISO 15489-1:2016. Information and documentation --- Records management --- Part 1: Concepts and principles}. Ginebra, Suiza. \\
    \textbf{Tipo de fuente} & Norma internacional \\
    \textbf{Identificador} & ISO 15489-1:2016 \\
    \textbf{Localización} & \texttt{\footnotesize iso.org --- catálogo de normas vigentes} \\
    \textbf{Constancia} & La norma aparece en el catálogo de la ISO y en los catálogos nacionales que la adoptan, entre ellos el Consejo Canadiense de Normas (SCC), AFNOR (Francia), SIS (Suecia) y DIN (Alemania), bajo el mismo número de referencia. \\
    \bottomrule
  \end{tabularx}
\end{table}

\begin{table}[H]
  \small
  \renewcommand{\arraystretch}{1.25}
  \begin{tabularx}{\textwidth}{@{}>{\raggedright\arraybackslash}p{3.0cm}>{\raggedright\arraybackslash}X@{}}
    \toprule
    \multicolumn{2}{@{}l}{\textbf{Ficha A-11} \quad\textnormal{\footnotesize clave de cita: \texttt{iso25010}}} \\
    \midrule
    \textbf{Referencia} & International Organization for Standardization (2011). \emph{ISO/IEC 25010:2011. Systems and software engineering --- SQuaRE --- System and software quality models}. Ginebra, Suiza. \\
    \textbf{Tipo de fuente} & Norma internacional \\
    \textbf{Identificador} & ISO/IEC 25010:2011 \\
    \textbf{Localización} & \texttt{\footnotesize iso.org/standard/35733.html} \\
    \textbf{Constancia} & Ficha oficial de la norma en el catálogo de la ISO, bajo el número de registro 35733. Aparece asimismo en los catálogos de EVS (Estonia), DIN (Alemania) y UNE (España). \\
    \bottomrule
  \end{tabularx}
\end{table}

\begin{table}[H]
  \small
  \renewcommand{\arraystretch}{1.25}
  \begin{tabularx}{\textwidth}{@{}>{\raggedright\arraybackslash}p{3.0cm}>{\raggedright\arraybackslash}X@{}}
    \toprule
    \multicolumn{2}{@{}l}{\textbf{Ficha A-12} \quad\textnormal{\footnotesize clave de cita: \texttt{iso29148}}} \\
    \midrule
    \textbf{Referencia} & International Organization for Standardization (2018). \emph{ISO/IEC/IEEE 29148:2018. Systems and software engineering --- Life cycle processes --- Requirements engineering}. Ginebra, Suiza. \\
    \textbf{Tipo de fuente} & Norma internacional \\
    \textbf{Identificador} & ISO/IEC/IEEE 29148:2018 \\
    \textbf{Localización} & \texttt{\footnotesize committee.iso.org/standard/72089.html} \\
    \textbf{Constancia} & Ficha oficial en el catálogo de la ISO bajo el número de registro 72089. La norma se encuentra también en la tienda de la IEC (publicación 64315) y en el catálogo del IEEE como 29148-2018. \\
    \bottomrule
  \end{tabularx}
\end{table}

\begin{table}[H]
  \small
  \renewcommand{\arraystretch}{1.25}
  \begin{tabularx}{\textwidth}{@{}>{\raggedright\arraybackslash}p{3.0cm}>{\raggedright\arraybackslash}X@{}}
    \toprule
    \multicolumn{2}{@{}l}{\textbf{Ficha A-13} \quad\textnormal{\footnotesize clave de cita: \texttt{owasp2021}}} \\
    \midrule
    \textbf{Referencia} & OWASP Foundation (2021). \emph{OWASP Top 10: The Ten Most Critical Web Application Security Risks}. \\
    \textbf{Tipo de fuente} & Documento técnico de referencia \\
    \textbf{Identificador} & Edición 2021 del OWASP Top 10 \\
    \textbf{Localización} & \texttt{\footnotesize owasp.org/Top10/} \\
    \textbf{Constancia} & Cada categoría de la edición 2021 tiene dirección propia y permanente dentro del sitio de la fundación; por ejemplo \texttt{owasp.org/Top10/2021/A01\_2021-Broken\_Access\_Control/} y la introducción en \texttt{A00\_2021\_Introduction/}. \\
    \bottomrule
  \end{tabularx}
\end{table}

\begin{table}[H]
  \small
  \renewcommand{\arraystretch}{1.25}
  \begin{tabularx}{\textwidth}{@{}>{\raggedright\arraybackslash}p{3.0cm}>{\raggedright\arraybackslash}X@{}}
    \toprule
    \multicolumn{2}{@{}l}{\textbf{Ficha A-14} \quad\textnormal{\footnotesize clave de cita: \texttt{scrumguide2020}}} \\
    \midrule
    \textbf{Referencia} & Schwaber, K. y Sutherland, J. (2020). \emph{The Scrum Guide: The Definitive Guide to Scrum. The Rules of the Game}. Noviembre de 2020. \\
    \textbf{Tipo de fuente} & Documento técnico de referencia \\
    \textbf{Identificador} & Edición de noviembre de 2020 \\
    \textbf{Localización} & \texttt{\footnotesize scrumguides.org/docs/scrumguide/v2020/\allowbreak 2020-Scrum-Guide-US.pdf} \\
    \textbf{Constancia} & El documento se distribuye en PDF desde el sitio de sus autores, con versiones paralelas en varios idiomas bajo la misma ruta \texttt{/v2020/}, entre ellas la edición en español latinoamericano. \\
    \bottomrule
  \end{tabularx}
\end{table}

\subsection*{Ingeniería de software y bases de datos}
\addcontentsline{toc}{subsection}{Ingeniería de software y bases de datos}

\begin{table}[H]
  \small
  \renewcommand{\arraystretch}{1.25}
  \begin{tabularx}{\textwidth}{@{}>{\raggedright\arraybackslash}p{3.0cm}>{\raggedright\arraybackslash}X@{}}
    \toprule
    \multicolumn{2}{@{}l}{\textbf{Ficha A-15} \quad\textnormal{\footnotesize clave de cita: \texttt{sommerville2011}}} \\
    \midrule
    \textbf{Referencia} & Sommerville, I. (2011). \emph{Software Engineering} (9.\textsuperscript{a} ed.). Boston: Addison-Wesley. \\
    \textbf{Tipo de fuente} & Libro \\
    \textbf{Identificador} & ISBN 978-0-13-703515-1 \\
    \textbf{Localización} & \texttt{\footnotesize Catálogo del editor y catálogos bibliotecarios} \\
    \textbf{Constancia} & La novena edición se registra con el ISBN indicado y figura en catálogos bibliotecarios universitarios y en el registro de ediciones de la obra, que documenta la serie completa desde la primera edición. \\
    \bottomrule
  \end{tabularx}
\end{table}

\begin{table}[H]
  \small
  \renewcommand{\arraystretch}{1.25}
  \begin{tabularx}{\textwidth}{@{}>{\raggedright\arraybackslash}p{3.0cm}>{\raggedright\arraybackslash}X@{}}
    \toprule
    \multicolumn{2}{@{}l}{\textbf{Ficha A-16} \quad\textnormal{\footnotesize clave de cita: \texttt{pressman2010}}} \\
    \midrule
    \textbf{Referencia} & Pressman, R. S. (2010). \emph{Software Engineering: A Practitioner's Approach} (7.\textsuperscript{a} ed.). Nueva York: McGraw-Hill. \\
    \textbf{Tipo de fuente} & Libro \\
    \textbf{Identificador} & ISBN 978-0-07-337597-7 \\
    \textbf{Localización} & \texttt{\footnotesize highered.mheducation.com/sites/0073375977/} \\
    \textbf{Constancia} & La séptima edición conserva centro de información propio en el sitio del editor McGraw-Hill, bajo el identificador 0073375977, y está catalogada en acervos universitarios con el ISBN indicado. \\
    \bottomrule
  \end{tabularx}
\end{table}

\begin{table}[H]
  \small
  \renewcommand{\arraystretch}{1.25}
  \begin{tabularx}{\textwidth}{@{}>{\raggedright\arraybackslash}p{3.0cm}>{\raggedright\arraybackslash}X@{}}
    \toprule
    \multicolumn{2}{@{}l}{\textbf{Ficha A-17} \quad\textnormal{\footnotesize clave de cita: \texttt{bass2021}}} \\
    \midrule
    \textbf{Referencia} & Bass, L., Clements, P. y Kazman, R. (2021). \emph{Software Architecture in Practice} (4.\textsuperscript{a} ed.). Boston: Addison-Wesley. \\
    \textbf{Tipo de fuente} & Libro \\
    \textbf{Identificador} & ISBN 978-0-13-688609-9 \\
    \textbf{Localización} & \texttt{\footnotesize Catálogo del editor y plataforma O'Reilly} \\
    \textbf{Constancia} & La cuarta edición se registra con el ISBN impreso indicado y con el ISBN electrónico 9780136885979, bajo el cual la obra se encuentra disponible íntegra, capítulo por capítulo, en la biblioteca digital de O'Reilly. \\
    \bottomrule
  \end{tabularx}
\end{table}

\begin{table}[H]
  \small
  \renewcommand{\arraystretch}{1.25}
  \begin{tabularx}{\textwidth}{@{}>{\raggedright\arraybackslash}p{3.0cm}>{\raggedright\arraybackslash}X@{}}
    \toprule
    \multicolumn{2}{@{}l}{\textbf{Ficha A-18} \quad\textnormal{\footnotesize clave de cita: \texttt{elmasri2016}}} \\
    \midrule
    \textbf{Referencia} & Elmasri, R. y Navathe, S. B. (2016). \emph{Fundamentals of Database Systems} (7.\textsuperscript{a} ed.). Boston: Pearson. \\
    \textbf{Tipo de fuente} & Libro \\
    \textbf{Identificador} & ISBN 978-0-13-397077-7 \\
    \textbf{Localización} & \texttt{\footnotesize pearson.com --- catálogo de la obra} \\
    \textbf{Constancia} & La séptima edición figura en el catálogo del editor Pearson con ficha propia, y está registrada en WorldCat con el número de control OCLC 913842106. \\
    \bottomrule
  \end{tabularx}
\end{table}

\begin{table}[H]
  \small
  \renewcommand{\arraystretch}{1.25}
  \begin{tabularx}{\textwidth}{@{}>{\raggedright\arraybackslash}p{3.0cm}>{\raggedright\arraybackslash}X@{}}
    \toprule
    \multicolumn{2}{@{}l}{\textbf{Ficha A-19} \quad\textnormal{\footnotesize clave de cita: \texttt{laraveldocs}}} \\
    \midrule
    \textbf{Referencia} & Laravel (2026). \emph{Laravel Documentation}. \\
    \textbf{Tipo de fuente} & Documentación técnica oficial \\
    \textbf{Identificador} & Dominio oficial del proyecto \texttt{laravel.com} \\
    \textbf{Localización} & \texttt{\footnotesize laravel.com/docs} \\
    \textbf{Constancia} & Documentación oficial del marco de trabajo, publicada por el propio proyecto y versionada por número de edición mayor dentro de la misma ruta. \\
    \bottomrule
  \end{tabularx}
\end{table}

\subsection*{Plataformas de gestión documental revisadas}
\addcontentsline{toc}{subsection}{Plataformas de gestión documental revisadas}

\begin{table}[H]
  \small
  \renewcommand{\arraystretch}{1.25}
  \begin{tabularx}{\textwidth}{@{}>{\raggedright\arraybackslash}p{3.0cm}>{\raggedright\arraybackslash}X@{}}
    \toprule
    \multicolumn{2}{@{}l}{\textbf{Ficha A-20} \quad\textnormal{\footnotesize clave de cita: \texttt{alfresco}}} \\
    \midrule
    \textbf{Referencia} & Hyland Software (2026). \emph{Alfresco Content Services: Enterprise Content Management}. \\
    \textbf{Tipo de fuente} & Sitio oficial del producto \\
    \textbf{Identificador} & Dominio oficial \texttt{alfresco.com} \\
    \textbf{Localización} & \texttt{\footnotesize alfresco.com} \\
    \textbf{Constancia} & Sitio oficial del producto, operado por Hyland Software tras la adquisición de Alfresco Software. La plataforma se revisó como antecedente en la sección 2.3 del presente trabajo. \\
    \bottomrule
  \end{tabularx}
\end{table}

\begin{table}[H]
  \small
  \renewcommand{\arraystretch}{1.25}
  \begin{tabularx}{\textwidth}{@{}>{\raggedright\arraybackslash}p{3.0cm}>{\raggedright\arraybackslash}X@{}}
    \toprule
    \multicolumn{2}{@{}l}{\textbf{Ficha A-21} \quad\textnormal{\footnotesize clave de cita: \texttt{openkm}}} \\
    \midrule
    \textbf{Referencia} & OpenKM (2026). \emph{OpenKM Document Management System}. \\
    \textbf{Tipo de fuente} & Sitio oficial del producto \\
    \textbf{Identificador} & Dominio oficial \texttt{openkm.com} \\
    \textbf{Localización} & \texttt{\footnotesize openkm.com} \\
    \textbf{Constancia} & Sitio oficial del producto. La plataforma aparece de manera recurrente en los recuentos especializados de sistemas de gestión documental de código abierto y se revisó como antecedente en la sección 2.3. \\
    \bottomrule
  \end{tabularx}
\end{table}

\begin{table}[H]
  \small
  \renewcommand{\arraystretch}{1.25}
  \begin{tabularx}{\textwidth}{@{}>{\raggedright\arraybackslash}p{3.0cm}>{\raggedright\arraybackslash}X@{}}
    \toprule
    \multicolumn{2}{@{}l}{\textbf{Ficha A-22} \quad\textnormal{\footnotesize clave de cita: \texttt{mayanedms}}} \\
    \midrule
    \textbf{Referencia} & Mayan EDMS (2026). \emph{Mayan EDMS: Free Open Source Electronic Document Management System}. \\
    \textbf{Tipo de fuente} & Sitio oficial del producto \\
    \textbf{Identificador} & Dominio oficial \texttt{mayan-edms.com} \\
    \textbf{Localización} & \texttt{\footnotesize mayan-edms.com} \\
    \textbf{Constancia} & Sitio oficial del proyecto, de código abierto. La plataforma aparece en los recuentos especializados de sistemas de gestión documental para entornos GNU/Linux y se revisó como antecedente en la sección 2.3. \\
    \bottomrule
  \end{tabularx}
\end{table}

\begin{table}[H]
  \small
  \renewcommand{\arraystretch}{1.25}
  \begin{tabularx}{\textwidth}{@{}>{\raggedright\arraybackslash}p{3.0cm}>{\raggedright\arraybackslash}X@{}}
    \toprule
    \multicolumn{2}{@{}l}{\textbf{Ficha A-23} \quad\textnormal{\footnotesize clave de cita: \texttt{seeddms}}} \\
    \midrule
    \textbf{Referencia} & SeedDMS (2026). \emph{SeedDMS: Free Document Management System}. \\
    \textbf{Tipo de fuente} & Sitio oficial del producto \\
    \textbf{Identificador} & Dominio oficial \texttt{seeddms.org} \\
    \textbf{Localización} & \texttt{\footnotesize seeddms.org} \\
    \textbf{Constancia} & Sitio oficial del proyecto, de código abierto y continuador de LetoDMS. La plataforma se revisó como antecedente en la sección 2.3. \\
    \bottomrule
  \end{tabularx}
\end{table}


\vspace{10pt}
\hrule
\vspace{8pt}

\section*{Nota sobre el alcance de la evidencia}

Las fichas acreditan la existencia y la localización de cada fuente, no el contenido que de ellas se cita; este último
queda respaldado por las citas puntuales del cuerpo del reporte. En el caso de las normas internacionales, el texto
íntegro es de acceso oneroso, por lo que la evidencia consignada es la ficha de catálogo del organismo emisor, que es
la constancia pública disponible. En el caso de los libros, la evidencia es el registro bibliográfico de la edición
citada, identificada por su ISBN.

\end{document}
